%%%%%%%%%%%%%%%%%%%%%%%%%%%%%%%%%%%%%%%%%%%%%%%%%%%%%%%%%%%%%%%%%%%%%%%%%%%%%%%%
%2345678901234567890123456789012345678901234567890123456789012345678901234567890
%        1         2         3         4         5         6         7         8

\documentclass[letterpaper, 10pt, journal]{IEEEtran}  % Comment this line out if you need a4paper

%\documentclass[a4paper, 10pt, conference]{ieeeconf}      % Use this line for a4 paper

% \IEEEoverridecommandlockouts                              % This command is only needed if 
                                                          % you want to use the \thanks command

% \overrideIEEEmargins                                      % Needed to meet printer requirements.

%In case you encounter the following error:
%Error 1010 The PDF file may be corrupt (unable to open PDF file) OR
%Error 1000 An error occurred while parsing a contents stream. Unable to analyze the PDF file.
%This is a known problem with pdfLaTeX conversion filter. The file cannot be opened with acrobat reader
%Please use one of the alternatives below to circumvent this error by uncommenting one or the other
%\pdfobjcompresslevel=0
%\pdfminorversion=4

% See the \addtolength command later in the file to balance the column lengths
% on the last page of the document

% The following packages can be found on http:\\www.ctan.org
%\usepackage{graphics} % for pdf, bitmapped graphics files
%\usepackage{epsfig} % for postscript graphics files
%\usepackage{mathptmx} % assumes new font selection scheme installed
%\usepackage{times} % assumes new font selection scheme installed
%\usepackage{amsmath} % assumes amsmath package installed
%\usepackage{amssymb}  % assumes amsmath package installed

% reference 
\bibliographystyle{IEEEtran}
\usepackage{graphicx}
\usepackage{caption}
\usepackage{subcaption}
\usepackage{booktabs}
\usepackage{array}
\usepackage{xcolor}
\usepackage{hyperref}
\hypersetup{%
  colorlinks=true,%
  linkcolor={red},
  citecolor={blue},
  urlcolor={magenta},
  bookmarksnumbered=true,%
  bookmarksopen=true}


\usepackage[normalem]{ulem} % for \uline (better-looking underline)
% If you already load ulem or dislike it, you can use \underline instead.
\usepackage{amsmath}
\usepackage{amsfonts}
\usepackage{multirow}
\newcommand{\todo}[1]{\textcolor{blue}{TODO: #1}}
\newcommand{\mllm}[1]{\textcolor{teal}{#1}}
\newcommand{\vlp}[1]{\textcolor{violet}{#1}}
\newcommand{\prompt}[1]{#1}
\newcommand{\qwentwob}{Qwen3-VL-2B}
\newcommand{\qweneightb}{Qwen3-VL-8B}
\newcommand{\qwensevenb}{Qwen2.5-VL-7B}


\title{\LARGE \bf

Decoding Pedestrian Crossing Intention from Egocentric Vision \\ 
via Vision Language Models
}


\author{Danya Li$^{1}$, Xiang Su$^{2}$, Yan Feng$^{3}$, and Rico Krueger $^{1}$% <-this % stops a space
% \thanks{*This work was not supported by any organization}% <-this % stops a space
\thanks{$^{1}$Danya Li and Rico Krueger are with Department of Technology, Management, and Economics, Technical University of Denmark, Denmark
        {\tt\small danli@dtu.dk, rickr@dtu.dk}}%
\thanks{$^{2}$Xiang Su is with Department of Computer Science and Department of Agricultural Sciences, University of Helsinki, Finland
    {\tt\small xiang.su@helsinki.fi}}
\thanks{$^{3}$Yan Feng is with Department of Transport and Planning, Delft University of Technology, Netherlands
    {\tt\small y.feng@tudelft.nl}
}
% \thanks{$^{2}$Bernard D. Researcheris with the Department of Electrical Engineering, Wright State University,
%         Dayton, OH 45435, USA
%         {\tt\small b.d.researcher@ieee.org}}%
}


\begin{document}



\maketitle
\thispagestyle{plain}
\pagestyle{plain}


%%%%%%%%%%%%%%%%%%%%%%%%%%%%%%%%%%%%%%%%%%%%%%%%%%%%%%%%%%%%%%%%%%%%%%%%%%%%%%%%
\begin{abstract}

Egocentric vision offers a first-person view of human perception and decision making, yet its potential for traffic-safety prediction remains underexplored. In this work, we study the decoding of pedestrian crossing intentions from short egocentric video clips. We approach this by formulating the task as a closed-ended visual question answering (VQA) problem and leveraging vision language models (VLMs) to predict the pedestrians' intent. We first benchmark three families of state-of-the-art VLMs in a zero-shot setting, finding that they achieve moderate gains over random guessing but exhibit limited higher-level traffic reasoning. Motivated by these findings, we further adapt VLMs to the target task using parameter-efficient fine-tuning. Our results show that the fine-tuned models substantially outperform their zero-shot counterparts and achieve a 9\% accuracy improvement over a specialized transformer-based baseline. Finally, we demonstrate that incorporating additional contextual cues, including ego motion, vehicle motion, and eye gaze, further improves predictive performance. In particular, the fine-tuned \qwentwob\ model guided by eye gaze and ego motion achieves a 14.5\% accuracy improvement over the transformer baseline, establishing a new state of the art for egocentric pedestrian intent decoding.
Code will be made available at \url{https://github.com/danyayay/EgoCross-VLM.git}.

\end{abstract}

\vspace{0.6em}
\noindent
\textbf{Keywords}: pedestrian intention prediction, vision--language models (VLMs), egocentric vision, eye gaze

%%%%%%%%%%%%%%%%%%%%%%%%%%%%%%%%%%%%%%%%%%%%%%%%%%%%%%%%%%%%%%%%%%%%%%%%%%%%%%%%
\section{INTRODUCTION}

% This template provides authors with most of the formatting specifications needed for preparing electronic versions of their papers. All standard paper components have been specified for three reasons: (1) ease of use when formatting individual papers, (2) automatic compliance to electronic requirements that facilitate the concurrent or later production of electronic products, and (3) conformity of style throughout a conference proceedings. Margins, column widths, line spacing, and type styles are built-in; examples of the type styles are provided throughout this document and are identified in italic type, within parentheses, following the example. Some components, such as multi-leveled equations, graphics, and tables are not prescribed, although the various table text styles are provided. The formatter will need to create these components, incorporating the applicable criteria that follow.

Ensuring pedestrian safety in increasingly complex urban environments is critically dependent on accurate anticipation of pedestrian behavior. Traditional approaches have predominantly relied on exocentric perspectives, such as vehicle-mounted or surveillance cameras, which offer a stable and global view of the scene \cite{zhang_pedestrian_2023, landry_predicting_2025}. 
However, these external viewpoints often fail to capture the pedestrian's first-person perception and fine-grained behavioral cues. In contrast, egocentric vision offers direct access to the pedestrian’s line of sight, providing rich information that is crucial to understand intention and action execution \cite{wang_egonav_2024, he_bridging_2026}. 

The growing interest in head-mounted display technologies (e.g., smart glasses) has created new opportunities for egocentric sensing. Although such devices are still in the early stages of adoption, their potential has spurred substantial research on egocentric vision indoors, enabling general-purpose egocentric assistants \cite{grauman_ego4d_2022, plizzari2024outlook, yang_egolife_2026}. Extending these capabilities to dynamic urban environments remains an important but underexplored challenge. Addressing this gap is central to proactive AI systems for pedestrian assistance, including navigation support \cite{wang_egonav_2024, qiu_egocognav_2025, pan2025lookout} and specialized aids for visually impaired users \cite{haghighi2024heads}. 

\begin{figure}[t]
    \centering
    \includegraphics[width=\linewidth]{figures/overview.png}
    \caption{A pipeline of our approach.}
    \label{fig:overview}
\end{figure}

Existing research on egocentric prediction has largely focused on motion forecasting. Prior work has studied egocentric trajectory prediction in crowded spaces \cite{qiu_egocentric_2022} and on sidewalks \cite{singh_krishnacam_2016, pan2025lookout}, typically using task-specific architectures such as convolutional neural networks \cite{singh_krishnacam_2016}, diffusion models \cite{wang_egonav_2024}, and transformer-based approaches \cite{qiu_egocognav_2025}. These methods often use heterogeneous inputs, ranging from monocular views \cite{pan2025lookout, qiu_egocognav_2025} and stereo images \cite{park2016egocentric} to panoramic videos \cite{wang_egonav_2024}. Such specialization can limit transferability and generalization to diverse, real-world urban scenarios.

To address these limitations, we explore the potential of pretrained vision--language models (VLMs) for egocentric intention prediction by reformulating the problem as a visual question answering (VQA) task.
We ask: \textbf{Can we infer pedestrians' future crossing intention from monocular first-person videos in urban environments?}
Specifically, we first evaluate the zero-shot reasoning capabilities of state-of-the-art VLMs to assess whether they already possess the perceptual and reasoning skills required for urban navigation. We then study the impact of parameter-efficient fine-tuning for adapting these models to traffic-specific prediction. Finally, by incorporating gaze signals and personal attributes, we investigate whether additional contextual cues can improve egocentric traffic reasoning beyond video-only perception.
An overview of our framework is presented in Fig.~\ref{fig:overview}.

The remainder of the paper is organized as follows: Sec.~\ref{sec:related_work} reviews related literature. Sec.~\ref{sec:methodology} and \ref{sec:experimental_setup} describe the methodology and experimental setup. Sec.~\ref{sec:results_n_analysis} present and analyze the results. We conclude the paper in Sec.~\ref{sec:conclusion} with limitations and future directions. 


\section{Related Work}
\label{sec:related_work}

\subsection{Pedestrian Intention Prediction}

Pedestrian intention prediction is a longstanding research topic in intelligent transportation and autonomous driving \cite{sharma2022pedestrian}. Most existing approaches operate under an \emph{exocentric} setting---from vehicle-mounted sensors, roadside cameras, or bird's-eye views---and infer what a pedestrian is about to do based on externally observable cues. 

In a typical pipeline, models combine multimodal inputs (e.g., tracked trajectories, body pose, and vehicle kinematics) with feature extractors and a downstream decoder/classifier. As summarized by recent surveys \cite{zhang_pedestrian_2023, landry_predicting_2025}, representative architectures span convolutional and recurrent models, graph-based interaction reasoning, and transformer-based sequence modeling.
More recently, VLMs have emerged as a promising paradigm for pedestrian-related traffic understanding. Several works probe multimodal large language models (MLLMs) in zero-shot settings to assess their traffic-scene understanding and reasoning \cite{huang_gpt-4v_2024, ham_omnipredict_2024, zambare2025seeing}. Other studies report gains via prompt engineering (e.g., hierarchical prompt templates \cite{azarmi_pedestrian_2025}), by using pretrained multimodal encoders as feature extractors \cite{uziel_optimizing_2025, munir_pedestrian_2025}, or by leveraging proprietary models for data annotation and teacher--student distillation \cite{gao_application_2025, ling2026vlmped}.

% Despite these advances, prior pedestrian intention prediction remains predominantly exocentric and models pedestrians as observed agents. This leaves an important gap: inferring intention from an \emph{egocentric} viewpoint that reflects the pedestrian's own perceptual evidence and decision making.

\subsection{Egocentric vision for behavior understanding}

% Egocentric vision studies videos captured from wearable sensors and cameras. 
The emergence of large-scale egocentric datasets has catalyzed research on general-purpose \emph{egocentric video understanding}. Many widely used datasets predominantly capture indoor daily activities, such as EPIC-KITCHENS \cite{damen_scaling_2018} featuring kitchen activities, Charades-Ego \cite{sigurdsson2018actor} consisting of everyday tasks, and Ego4D \cite{grauman_ego4d_2022} with more diverse first-person videos, still largely centered on household activities. These resources have enabled research problems including human--object interaction understanding, action recognition, and anticipatory/predictive modeling.
Video--language pretraining (VLP) approaches aim to learn aligned representations that better account for the characteristics of first-person videos (e.g., frequent viewpoint changes and strong ego motion). For example, EgoVLPv2 learns cross-modal representations from large-scale egocentric video--text data \cite{pramanick_egovlpv2_2023}, and recent efforts explore building foundation models tailored to egocentric video \cite{pei2024egovideo}. 
% In parallel, the community has begun to emphasize long-term egocentric understanding with memory-like capabilities (e.g., habit tracking and task support) \cite{yang_egolife_2026}.

% More recently, vision--language models (VLMs) and video--language pretraining have been extended to egocentric videos, enabling language-conditioned understanding beyond fixed label spaces. For example, EgoVLPv2 learns aligned representations from large-scale first-person video and text data \cite{pramanick_egovlpv2_2023}, and multimodal LLMs have demonstrated growing capability in open-ended egocentric reasoning \cite{patel_advancing_2025, peng2025eye}. Despite this progress, most egocentric VLM benchmarks and training data remain biased toward indoor activity understanding, whereas outdoor traffic interactions demand distinct reasoning about moving vehicles, right-of-way negotiation, and short-horizon action selection.

Despite rapid progress on indoor activity understanding, egocentric \emph{urban} understanding with highly dynamic, safety-critical interactions remains comparatively underexplored.
One relevant line of work is \emph{egocentric motion prediction} using deep learning, which forecasts the future movement of the camera wearer. Early approaches explored future localization from egocentric images or stereo inputs, for example via nearest-neighbor retrieval from a single egocentric image \cite{singh_krishnacam_2016} or stereo-based EgoRetinal representations \cite{park2016egocentric}. More recent methods predict wearer trajectories and head motion in crowded and urban environments \cite{qiu_egocentric_2022, pan2025lookout}, including diffusion-based modeling of future trajectories from panoramic videos \cite{wang_egonav_2024} and transformer-style sequence modeling with cognition-aware uncertainty \cite{qiu_egocognav_2025}. Wearable-camera datasets developed for assistive mobility further highlight the value of pedestrian-centric sensing for motion forecasting \cite{haghighi2024heads}. In our work, we move beyond motion forecasting and study short-horizon egocentric \emph{intention} decoding in traffic interactions using VLMs.

\subsection{Video question answering}


Video question answering (VideoQA) is a central task in video-language understanding, where models answer natural-language questions about dynamic visual content. Existing VideoQA tasks are commonly categorized into factoid and inference questions \cite{zhong2022video}. Factoid questions focus on directly observable evidence, such as locations, objects, attributes, actions, and activities, and mainly evaluate question understanding and visual recognition \cite{zhong2022video,xu2017video,jang2017tgif,yu2019activitynet}.

Recent benchmarks increasingly emphasize inference-oriented VideoQA, which requires reasoning over relations among objects, actions, events, states, and agents. These benchmarks evaluate temporal reasoning, such as action order, repetition, and state transitions \cite{jang2017tgif,xiao2021next}; spatial and compositional spatio-temporal reasoning over object-action relations and event structures \cite{jang2017tgif,grunde2021agqa}; causal reasoning through explanatory and predictive questions \cite{xiao2021next,li2022from}; and counterfactual reasoning through hypothetical ``what-if'' questions \cite{li2022from,jia2022egotaskqa}.

In egocentric settings, VideoQA further supports the evaluation of goal-directed task understanding, including action dependencies, action effects, intents, goals, and agents' beliefs \cite{jia2022egotaskqa}. Recent work also studies intent-oriented VideoQA in both third-person and egocentric videos, including context-aware intent reasoning and gaze-guided egocentric intent understanding \cite{li2023intentqa,peng2025eye}.




\section{Methodology}
\label{sec:methodology}

\subsection{Problem definition}

We study pedestrian crossing intention prediction from egocentric observations. Given a 2-second egocentric video clip $V$ and optional contextual information $C$, the goal is to predict whether the pedestrian will \emph{cross} or \emph{yield} within a 1-second future horizon. Formally, the task is defined as
\begin{equation}
    y = \mathcal{F}_{\Theta}(V, C),
\end{equation}
where $V \in \mathbb{R}^{T \times H \times W \times 3}$ denotes the input video with $T$ frames of spatial resolution $H \times W$, $C$ denotes additional contextual information described in Sec.~\ref{sec:contex}, and $y \in \mathcal{Y}=\{\text{cross}, \text{yield}\}$ denotes the predicted crossing intention.

We formulate intention prediction as a closed-ended visual question answering (VQA) task in order to leverage the visual understanding, language grounding, and reasoning priors of large pretrained VLMs. Specifically, the model receives the video $V$, additional contextual information $C$, and an intention query $Q$, and selects an answer from the predefined answer set $\mathcal{Y}$:
\begin{equation}
    A = \text{argmax}_{a \in \mathcal{Y}} p_{\Theta}(a \mid V, C, Q),
\end{equation}
where $A$ is the predicted answer and $p_{\Theta}$ denotes the answer probability estimated by the VLM.


\subsection{Model selection}
We consider two complementary families of video-based VLMs. The first family comprises \textbf{MLLMs} that directly process visual (video) inputs and are pretrained on large and diverse corpora, yielding broad world knowledge and general reasoning ability.
To balance performance and computational efficiency, we adopt Qwen3-VL-8B-Instruct and its lightweight variant Qwen3-VL-2B-Instruct \cite{qwen3technicalreport}, as well as Qwen2.5-VL-7B-Instruct, which has shown strong performance on egocentric understanding benchmarks \cite{peng2025eye}. We also evaluate InternVL3 \cite{zhu2025internvl3} using InternVL3-2B and InternVL3-8B. For the 7B and 8B models, we consider both 16-bit and 8-bit quantization.

The second family comprises \textbf{VLPs}, a de facto paradigm for video - text tasks. These models explicitly pretrain cross-modal representations on egocentric video and question answering data to improve video--text alignment \cite{pramanick_egovlpv2_2023}. However, their training data is often dominated by indoor activities \cite{damen_scaling_2018, grauman_ego4d_2022}. We select GroundVQA \cite{di_grounded_2024} (hereafter \textit{VLP}), a state-of-the-art architecture with dual video and language encoders and a cross-modal fusion module.


\subsection{Prompt design for VLMs}
We evaluate three prompting strategies, namely, standard prompting, text Chain-of-Thought (CoT) prompting \cite{wei_chain--thought_2022}, and visual prompting, because they probe complementary model behaviors that are all relevant to traffic-scene intention decoding. Standard prompting measures the model's out-of-the-box ability to map egocentric perception to a discrete action from minimal instructions, and therefore serves as our primary zero-shot reference. Text CoT prompting tests whether eliciting explicit intermediate reasoning in the language domain improves decision making beyond direct classification. Visual prompting, in contrast, injects task-relevant visual cues directly into the input to encourage spatial grounding and temporal consistency without modifying model parameters \cite{yang2023set,yang2023fine}.

For \textbf{standard prompt}, to enforce a structured response and fairly assess model capabilities, we design a constrained prompt: \prompt{``What is your most likely action in the next 1 second based on what you saw in the egocentric video for the past 2 seconds? Choose one option: (A) cross (B) yield.''}

For \textbf{text CoT prompting}, we consider a simple and an advanced variant.
First, we use \prompt{``Let's think step by step''} \cite{kojima_large_2023} to activate the chain-of-thought process. 
To further encourage multi-step reasoning, we use \prompt{``Analyze the egocentric video. First, describe the visual elements related to the crossing task. Second, evaluate attention presence, perceived proximity, and perceived risk. Third, explain the logic connecting these elements. Finally, provide the final answer. Output format: Reasoning: [maximum 5 sentences about your reasoning]. Answer: [just the letter and option].''}

To construct the \textbf{visual prompt}, we augment the video input with offline visual cues, inspired by Set-of-Mark prompting for visual grounding \cite{yang2023set}. Specifically, we use the open-vocabulary detector GroundingDINO \cite{liu2024grounding} to detect task-relevant objects; in our setting, we focus on the automated vehicle and the white crossing-circle goal. The detections are refined via geometry-based filtering, cross-label suppression, and lightweight temporal tracking based on SORT \cite{bewley2016simple}, ensuring stable identifiers across sampled frames. We then render the processed detections as set-of-marks overlays, where tracked objects are annotated with persistent numeric IDs. Examples are shown in Fig.~\ref{fig:visual_prompt}. 

\begin{figure*}[t]
\centering
    \begin{subfigure}{\textwidth}
    \centering
    \includegraphics[width=\textwidth]{figures/P19S6-3__f8_fps4_dur2__som_overlay_groundingdinotiny_011fc2e60c.pdf}
    \caption{Example: Consistent tracking of the same objects over time.}
    \label{fig:stable_version}
    \end{subfigure}
    
    \begin{subfigure}{\textwidth}
    \centering
    \includegraphics[width=\textwidth]{figures/P23S6-14__f8_fps4_dur2__som_overlay_groundingdinobase_011fc2e60c.pdf}
    \caption{Example: Identification of new objects when turning head.}
    \label{fig:turning_version}
    \end{subfigure}

    \caption{Examples of visual prompt overlay on raw videos. }
    \label{fig:visual_prompt}
\end{figure*}

For video encoding, we consider two input representations. The first uses direct video inputs when the model natively supports video encoding. For the Qwen model family, the vision--language architecture employs multimodal rotary position embeddings to split videos into 3D chunks and to encode absolute temporal indices in the positional embedding \cite{qwen3technicalreport}. The second representation uses interleaved inputs, where each video is treated as an ordered sequence of individual frames and timestamps and image frames are alternated. 

\subsection{Representing contextual information}
\label{sec:contex}


We consider additional contextual information beyond the video input, including ego motion, vehicle motion, eye gaze, and personal attributes \cite{li2026eye, li2025analyzing}, all of which have been shown to be relevant to pedestrian intention prediction. \emph{Ego motion} describes the pedestrian's current position and speed, while \emph{vehicle motion} describes the vehicle's current position and speed. We also incorporate fine-grained eye-gaze signals to explicitly represent the pedestrian's line of sight \cite{li_challenges_2025}. Following \cite{peng2025eye, li2026eye}, we design two gaze-guided representations: (1) \emph{gaze direction}, where textual prompts describe the gaze orientation in the top-down view represented in degrees, indicating the direction to which the pedestrian is attending; and (2) \emph{gaze on screen}, where normalized gaze fixation coordinates on the image plane are represented as ratios in $[0,1]$. Personal attributes are incorporated by including demographic information in the prompt. For example: \prompt{``You are a 27-year-old female with a Master's degree or equivalent education level. Your dominant hand in daily life is your left hand. You walk every day...''.} In addition, we consider a visual gaze representation by rendering gaze fixations as red dots directly on the raw video frames. No overlay is rendered during saccades.

For time-varying contextual information, we further examine how it should be integrated into the prompt. Specifically, we consider two strategies. The first strategy directly appends the contextual information to the end of the standard prompt, under the assumption that the temporal trend can be inferred when the context is provided in a compact form; we refer to this strategy as \emph{preface}. The second strategy inserts contextual information alongside the corresponding timestamps when interleaved inputs are used. This design assumes that temporal alignment between contextual cues and video frames may help the model associate the context with specific visual evidence in the video; we refer to this strategy as \emph{interleaved}.



\subsection{Fine-tuning strategies}
We employ supervised fine-tuning (SFT) on annotated triplets of the form $(V, C, Q, A)$ to adapt pretrained models to egocentric traffic reasoning. 
To maintain computational efficiency while preserving pretrained knowledge, we adopt Low-Rank Adaptation (LoRA) \cite{hu_lora_2021}. LoRA freezes the original model parameters $\Theta_0$ and introduces trainable low-rank updates into selected transformer layers. For a weight matrix $W_0 \in \mathbb{R}^{d \times k}$, the adapted weight is given by
\begin{equation}
    W = W_0 + \Delta W, \qquad \Delta W = BA,
\end{equation}
where $B \in \mathbb{R}^{d \times r}$, $A \in \mathbb{R}^{r \times k}$, and $r \ll \min(d,k)$. During training, only the low-rank parameters $A$ and $B$ are updated.

The model is optimized using the standard cross-entropy loss over the closed answer set:
\begin{equation}
    \mathcal{L}
    = - \frac{1}{N} \sum_{i=1}^{N}
    \log p_{\Theta}(a_i \mid V_i, C_i, Q_i),
\end{equation}
where $a_i$ is the ground-truth answer for the $i$-th training example.

To analyze which components benefit most from domain adaptation, we selectively apply LoRA to different parts of the VLM: (1) the language encoders, (2) the cross-modal fusion modules, and (3) both the language and cross-modal components jointly. The vision encoder is kept frozen, following its role as a pretrained visual feature extractor in VLPs \cite{di_grounded_2024}.


\section{Experimental Setup}
\label{sec:experimental_setup}

Due to the limited availability of real-world egocentric data for pedestrian intention prediction in urban environments, we utilize a VR-based dataset from \cite{li2026eye}. This dataset captures egocentric pedestrian navigation alongside automated shuttles in a shared space. It provides synchronized egocentric videos, eye gaze tracking, demographic profiles, and movement trajectories.


To ensure data quality, we retain only critical interactions occurring prior to the intersection crossing point.
The data is segmented into 2-second observation windows and 1-second prediction horizons with a 0.5-second stride, yielding 6,047 QA samples.
To prevent data leakage and evaluate robust generalization, we partition the dataset at the participant-level into training, validation, and test sets using a 6:1:3 split.

\textbf{Labeling} 
We define the binary intention labels—``cross'' and ``yield''—based on pedestrian kinematic behavior within the 1-second future horizon.
Labels are determined by the duration for which the pedestrian speed exceeds a predefined crossing threshold \cite{li2026eye}.
Specifically, if the speed remains above the threshold for the majority of the horizon, the sample is labeled ``cross''; otherwise, it is labeled ``yield''. This results in 2,486 crossing and 3,561 yielding samples. To mitigate the resulting mild class imbalance, we apply random under-sampling to the majority class during subsequent training. 

\textbf{Baselines}
We establish a video-only deep learning baseline to forecast crossing intention.
Frame-level visual features are extracted using a pretrained CLIP backbone \cite{radford_learning_2021}, and the resulting feature sequence is processed by a transformer encoder followed by a classification head. This baseline is denoted as \textit{CLIP+Transformer} hereafter. 
% We validated the backbone and main architecture combination and found out the best baseline to achieve is a CLIP-based backbone with a LSTM-based architecture. 
The model is trained for up to 100 epochs using the Adam optimizer (learning rate 0.001, batch size 64) with early stopping to prevent overfitting. In addition, we consider two other simple baselines: (1) always predicting the majority class, and (2) random guessing. All baseline results are reported in Tab.~\ref{tab:zeroshot_egocentric}.

\textbf{Fine-tuning}
\todo{}

\textbf{Metrics}
We evaluate performance using accuracy, following standard practice in close-ended question answering (CloseQA) tasks \cite{patel_advancing_2025, peng2025eye}. To account for class imbalance, we additionally report the macro F1 score.
To mitigate choice-order bias, we randomize the associate between answer option and intention labels. Furthermore, we set the sampling temperature to 0 for deterministic generation, thereby isolating the model reasoning capabilities from stochastic variation.

\textbf{Generalization checks} 
To examine whether the fine-tuned model genuinely used visual and temporal evidence rather than exploiting textual or class-label priors, we conducted a counterfactual input ablation study on the complete test set of 1,752 samples. The original evaluation used eight uniformly sampled frames presented in chronological order with interleaved timestamps. We then evaluated the same fine-tuned model under six input interventions while keeping the questions, prompts, preprocessing parameters, and deterministic decoding settings unchanged. In the no-video condition, only the textual prompt was provided. In the black-frame and random-noise conditions, the sampled frames were replaced with visually uninformative inputs of identical dimensions. In the mismatched-video condition, each question was paired with frames from a different test video, thereby preserving the visual-input format while removing correspondence between the question and the observed scene. To assess temporal sensitivity, we additionally randomized or reversed the order of the sampled frames while retaining their original temporal input positions. Frame shuffling was repeated using three fixed random seeds (42--44), and its results are reported as the mean and sample standard deviation across runs. These interventions were designed to distinguish reliance on language priors, input-specific visual evidence, and correctly ordered temporal dynamics.

To verify that fine-tuning on our domain-specific crossing-intention VQA task does not come at the cost of the model's general visual reasoning capabilities, we additionally evaluate our fine-tuned models on MMBench (English dev split)~\cite{duan2024vlmevalkit}, a standard multimodal benchmark covering everyday visual question answering across perception, reasoning, and knowledge categories unrelated to our task. For each fine-tuning configuration, we draw a stratified random subsample of 100 questions per category from MMBench's category taxonomy and evaluate both the base Qwen3-VL model and its LoRA fine-tuned counterpart on the identical set of questions under deterministic (greedy) decoding, repeating this procedure across three random seeds to account for sampling variance. We report overall and per-category multiple-choice accuracy (mean and standard deviation across seeds) for both models and compare the resulting deltas to determine whether fine-tuning on our narrow VQA task induces measurable degradation in the model's broader conversational and reasoning ability, or whether performance on general-purpose questions is preserved.


\section{Results and Analyses}
\label{sec:results_n_analysis}

\subsection{Zero-shot performance using solely egocentric videos}

We first benchmark the zero-shot capability of pretrained MLLMs and a VLP model against standard baselines (Tab.~\ref{tab:zeroshot_egocentric}). With the standard prompt, none of the evaluated VLMs surpass the CLIP+Transformer baseline (Acc.=0.727, M-F1=0.724). Nonetheless, several MLLMs perform above random guessing, suggesting that they capture task-relevant semantics from egocentric traffic videos. In contrast, the VLP model exhibits a strong class bias, predicting ``yield'' almost exclusively, which yields majority-level accuracy (0.567) but low macro F1 (0.378).

Adding text CoT prompting to Qwen does not improve performance, while substantially increasing latency (from 0.1--0.3~seconds to roughly 2--4~seconds per sample; Tab.~\ref{tab:zeroshot_egocentric}). Moreover, the more elaborate multi-step CoT prompt increases the tendency to default to ``yield''. This deterioration observation aligns with recent evidence that explicit CoT may degrade visual-spatial reasoning in multimodal models \cite{dai2026latentomni}. A plausible explanation is that text CoT prompting amplifies linguistic priors while weakening reliance on fine-grained visual evidence.

Motivated by this observation, we further examine visual prompting via set-of-marks overlays. However, this strategy does not yield a measurable improvement over standard prompting. This may be attributed to: (i) imperfect detection and tracking that introduce inconsistent identities, (ii) limited visual diversity in the VR environment, and (iii) a suboptimal choice of prompted object categories for intention decoding.

\begin{table*}[t]
\small
\centering
\caption{Baselines, zero-shot VLMs, and fine-tuned models using only egocentric video. \textbf{Bold} indicates the best result within each group; \dag~denotes 8-bit quantization. ``Time'' reports per-sample inference latency in seconds.}
\label{tab:zeroshot_egocentric}
% \resizebox{\linewidth}{!}{%
\begin{tabular}{lll!{\vrule}ccc}
\toprule
\textbf{Group} & \textbf{Model} & \textbf{Prompt} & \textbf{Acc.} & \textbf{M-F1} & \textbf{Time (s)} \\
\cmidrule(r){1-3} \cmidrule(l){4-6}
Baselines & Majority & --- & 0.567 & 0.362 & --- \\
          & Random   & --- & 0.500 & 0.497 & --- \\ 
          & CLIP+Transformer & --- & \textbf{0.727} & \textbf{0.724} & --- \\
\midrule
Zero-shot   & Qwen3-VL-2B & Standard     & 0.593 & 0.580 & 0.12 \\
            & Qwen3-VL-8B\dag & Standard & 0.579 & 0.559 & 0.32 \\ 
            & Qwen2.5-VL-7B\dag & Standard & \textbf{0.632} & \textbf{0.579} & 0.33 \\
            & Intern3VL-2B & Standard & 0.586 & 0.446 & 0.72 \\ 
            & Intern3VL-8B\dag & Standard & 0.594 & 0.475 & 1.35 \\ 
            & VLP  & Standard     & 0.567 & 0.378 & 0.01 \\
            
            \cmidrule(r){2-3} \cmidrule(l){4-6}
            & \qwensevenb\dag & CoT (simple) & 0.568 & 0.543 & 2.32 \\
            & \qwensevenb\dag & CoT (multi)  & 0.594 & 0.371 & 3.78 \\ 

            \cmidrule(r){2-3} \cmidrule(l){4-6}
            & \qwensevenb\dag & Visual prompt & 0.605 & 0.526 & 0.33 \\ 
            
\midrule
Fine-tuned & \qwentwob & Standard    & 0.777 & 0.773 & 0.12 \\ 
           & VLP  & Standard    & \textbf{0.789} & \textbf{0.787} & 0.007 \\
\bottomrule
\end{tabular}
% }
\end{table*}

We provide qualitative CoT examples in Fig.~\ref{fig:reasoning}. Although the model often identifies salient entities, it can still misread dynamic states. For example, in Fig.~\ref{fig:bad_ex}, the model recognizes the vehicle and the pedestrian's goal but incorrectly infers that the shuttle is moving and has already passed; in reality, the automated shuttle has stopped in front of the pedestrian.



\begin{figure}[t] % 't' suggests top placement
    \centering
    \begin{subfigure}{0.95\linewidth}
        \centering
        \includegraphics[width=\linewidth]{figures/ablation_interleaved_inputs.png}
        \caption{Effect of interleaved inputs.}
        \label{fig:ablation_zeroshot_interleaved_inputs}
    \end{subfigure}
    \hfill % Adds spacing between the two images
    \begin{subfigure}{0.95\linewidth}
        \centering
        \includegraphics[width=\linewidth]{figures/ablation_video_frame_number.png}
        \caption{Effect of sampled video frame number.}
        \label{fig:ablation_zeroshot_video_frame_number}
    \end{subfigure}
    \caption{Zero-shot ablations on temporal input representation. (a) Effect of representing the clip as interleaved timestamps and frames. (b) Effect of the number of sampled video frames.}
    \label{fig:ablation_zeroshot}
\end{figure}

To further analyze zero-shot performance, we conduct an ablation study on the effects of video representation strategy and video frame sampling rate. The results are shown in Fig.~\ref{fig:ablation_zeroshot}.

\paragraph{Effect of video representation strategy} Fig.~\ref{fig:ablation_zeroshot_interleaved_inputs} shows that interleaving timestamps and frames generally outperforms non-interleaved inputs for the Qwen family. This resonates with recent findings that a simple timestamp--frame interleaving can improve temporal reasoning across models on localization-style tasks \cite{Meinardus_2025_ICCV}.

\paragraph{Effect of the sampled video frame numbers}

We assess how video sampling rate affects zero-shot performance. Fig.~\ref{fig:ablation_zeroshot_video_frame_number} shows two trends: for larger models (e.g., \qweneightb), increasing the number of frames from 4 to 32 generally degrades performance, whereas for the smaller model (\qwentwob), additional frames are beneficial. This suggests that larger models may be more sensitive to redundant or noisy inputs, while smaller models may benefit from additional visual evidence to compensate for limited capacity.



Overall, the zero-shot results indicate that models pretrained on broad web-scale data or predominantly indoor egocentric corpora can extract some decision-relevant cues in traffic-like interactions, but still fall short of strong task performance without adaptation.




\subsection{Fine-tuning performance using solely egocentric videos}

\begin{figure}[t] % 't' suggests top placement
    \centering
    \begin{subfigure}{0.48\linewidth}
        \centering
        \includegraphics[width=\linewidth]{figures/qwen_finetuning.png}
        \caption{Qwen fine-tuning result.}
    \end{subfigure}
    \hfill % Adds spacing between the two images
    \begin{subfigure}{0.48\linewidth}
        \centering
        \includegraphics[width=\linewidth]{figures/vlp_finetuning.png}
        \caption{VLP fine-tuning result.}
    \end{subfigure}
    \caption{Comparison of parameter-efficient fine-tuning strategies for both VLMs. ``lan.'' tunes language modules, ``cm.'' tunes cross-modal modules, and ``lan.+cm.'' tunes both.}
    \label{fig:rq2_egocentric_only_ft}
\end{figure}

\begin{figure*}[p]
\centering
    \begin{subfigure}{\textwidth}
    \centering
    \includegraphics[width=\textwidth]{figures/badex_p30s3-17_3pics.pdf}
    \caption{Failure case: Misinterpretation of vehicle dynamics. Both models fail to recognize the automated shuttle has stopped for the pedestrian. Without a clear understanding of the vehicle’s intent, the models incorrectly estimate the safety of the path.}
    \label{fig:bad_ex}
    \end{subfigure}
    
    \begin{subfigure}{\textwidth}
    \centering
    \includegraphics[width=\textwidth]{figures/goodex_p48s6-1_3pics.pdf}
    \caption{Success case: Gaze-informed hazard perception. The gaze-overlaid model correctly anticipates potential hazards. This shows gaze signals helps mimic human-like situational awareness. }
    \label{fig:good_ex}
    \end{subfigure}

    \caption{Qualitative analysis of Qwen outputs under CoT prompting. We contrast egocentric-only input (left) with gaze-augmented input (right). Highlight colors denote reasoning quality: red (incorrect), green (correct), yellow (partially correct), and blue (context-relevant).}
    \label{fig:reasoning}
\end{figure*}


Building on the above findings, we apply parameter-efficient fine-tuning to adapt the pretrained VLMs to the target task while preserving their general knowledge. After adaptation, both models outperform the video-only CLIP+Transformer baseline (Tab.~\ref{tab:zeroshot_egocentric}): Qwen reaches Acc.=0.777, while VLP reaches Acc.=0.789. These improvements suggest that language-centric pretraining provides useful priors for mapping egocentric visual observations to discrete crossing decisions, but that domain adaptation remains necessary to capture traffic-specific interactions.




The optimal fine-tuning strategy varies across architectures (Fig.~\ref{fig:rq2_egocentric_only_ft}). For Qwen, updating only the cross-modal modules is sufficient to achieve strong performance, whereas VLP performs best when both the language and cross-modal components are adapted. Notably, VLP exhibits the larger improvement, suggesting greater adaptability to the target setting, potentially due to its egocentric pretraining.





\subsection{Impact of contextual information in zero-shot settings}

We next evaluate whether additional contextual cues improve zero-shot intention decoding. For the VLP model, adding contextual information leads to hallucinated outputs in approximately 5\% of cases; we therefore exclude it from this comparison. Tab.~\ref{tab:zeroshot_context_guided_summary} reports the results for \qwensevenb\ under different context configurations.

Three main observations emerge from the results. 
First, gaze-overlaid videos generally match or outperform raw videos across most context configurations, as indicated by the underlined entries in Tab.~\ref{tab:zeroshot_context_guided_summary}. Even without textual context, gaze overlay slightly improves accuracy over raw video input (0.640 vs. 0.632).
Second, adding a single contextual cue, such as gaze direction, gaze-on-screen coordinates, or vehicle motion, can degrade performance relative to the no-context setting. In contrast, combining these cues with ego motion yields more consistent gains. For example, with gaze-overlaid inputs, adding \emph{gaze on screen} together with \emph{ego motion} improves accuracy from 0.640 to 0.663. This suggests that ego-motion conditioning helps contextualize additional signals and stabilizes intention inference.
Third, \emph{ego motion} alone is already informative in the zero-shot setting, indicating that some contextual modalities may not be readily interpretable to the model without domain-specific training.


% \begin{table}[t]
% \centering
% \caption{Zero-shot performance with context guidance. For each setting, we report the best result over the context-FPS sweep. \textbf{Bold} number denotes the best performance in the group.}
% \label{tab:zeroshot_context_guided_summary}
% \begin{tabular}{ll|cc}
% \toprule
% Video Input & Context & Acc. & M-F1 \\
% \midrule
% Raw videos & - & \textbf{0.632} & \textbf{0.579} \\
%     & Gaze direction & 0.590 & 0.572 \\ 
%     & Gaze on screen & 0.592 & \textbf{0.579} \\
%     & Vehicle motion & 0.628 & 0.567 \\ 
%     \cmidrule(r){2-2} \cmidrule(l){3-4}
%     & Ego motion & 0.653 & \textbf{0.641} \\
%     & Ego motion, Gaze direction & \textbf{0.654} & 0.628 \\
%     & Ego motion, Gaze on screen & 0.646 & 0.636 \\
%     & Ego motion, Vehicle motion & 0.652 & 0.634 \\
% \midrule

% Gaze overlay & - & \textbf{0.640} & 0.574 \\
%     & Gaze direction & 0.622 & \textbf{0.592} \\ 
%     & Gaze on screen & 0.619 & 0.590 \\
%     & Vehicle motion & 0.636 & 0.564 \\ 
%     \cmidrule(r){2-2} \cmidrule(l){3-4}
%     & Ego motion & 0.640  & 0.588 \\
%     & Ego motion, Gaze direction & 0.657 & 0.614 \\
%     & Ego motion, Gaze on screen & \textbf{0.663} & \textbf{0.646} \\
%     & Ego motion, Vehicle motion & \textbf{0.663} & 0.637 \\
% \bottomrule
% \end{tabular}
% \end{table}

\begin{table}[t]
\small
\centering
\caption{Zero-shot performance with contextual guidance. For each context configuration, we report the best result over the context sampling-rate sweep. \uline{Underline} marks the better video input (gaze overlay vs. raw) within the same context and metric; \textbf{bold} marks the best context overall for each metric.}
\label{tab:zeroshot_context_guided_summary}
\begin{tabular}{l|cc|cc}
\toprule
& \multicolumn{2}{c|}{Raw videos} & \multicolumn{2}{c}{Gaze overlay} \\
\cmidrule(l){2-3} \cmidrule(l){4-5}
Context & Acc. & M-F1 & Acc. & M-F1 \\
\midrule
-- & 0.632 & \uline{0.579} & \uline{0.640} & 0.574 \\
Gaze direction & 0.590 & 0.572 & \uline{0.622} & \uline{0.592} \\
Gaze on screen & 0.592 & 0.579 & \uline{0.619} & \uline{0.590} \\
Vehicle motion & 0.628 & \uline{0.567} & \uline{0.636} & 0.564 \\
\cmidrule(r){1-1} \cmidrule(l){2-3} \cmidrule(l){4-5}
Ego motion & 0.653 & 0.641 & \uline{\textbf{0.664}} & \uline{0.643} \\
Ego motion, gaze direction & 0.654 & \uline{0.628} & \uline{0.657} & 0.614 \\
Ego motion, gaze on screen & 0.646 & 0.636 & \uline{0.663} & \uline{\textbf{0.646}} \\
Ego motion, vehicle motion & 0.652 & 0.634 & \uline{0.663} & \uline{0.637} \\
\bottomrule
\end{tabular}
\end{table}



Qualitative results (Fig.~\ref{fig:reasoning}) suggest that gaze overlay may enhance situational awareness, but in an implicit way. The generated reasoning rarely references gaze explicitly; nevertheless, as in Fig.~\ref{fig:good_ex}, the gaze-overlaid input can lead to more accurate projection of likely near-future events (e.g., inferring that an interaction with other agents may be ongoing, making it prudent to wait). Overall, these findings indicate that current pretrained models still struggle to reliably exploit context via prompting alone.

To further analyze these effects, we conduct an ablation study on context formatting and textual sampling rate when contextual cues are provided.

\paragraph{Effect of context formatting}
Tab.~\ref{tab:zeroshot_context_preface_interleaved} compares two formatting strategies: (i) appending contextual information as a preface and (ii) interleaving contextual information with video frames. The preface strategy consistently yields better results, whereas interleaving can substantially degrade performance. We speculate that interleaved text may disrupt the model's temporal visual processing, while a preface allows the model to condition globally on the contextual information before interpreting the frame sequence.

\begin{table}[t]
\small
\centering
\caption{Zero-shot performance with contextual guidance (gaze-overlay video input), comparing two context formats: preface vs. interleaved. Both use identical inputs (same text sampling rate). \uline{Underline} marks the better format; \textbf{bold} denotes the best context overall for each metric.}
\label{tab:zeroshot_context_preface_interleaved}
\begin{tabular}{ll|cc}
\toprule
Context & Text format? & Acc. & M-F1 \\
\midrule
Ego motion & Preface & \uline{\textbf{0.664}}  & \uline{\textbf{0.643}} \\
 & Interleaved & 0.662  & 0.624 \\
\cmidrule(r){1-2} \cmidrule(l){3-4}
Ego motion, gaze direction & Preface & \uline{0.623} & \uline{0.538} \\
 & Interleaved & 0.618 & 0.532 \\
\cmidrule(r){1-2} \cmidrule(l){3-4}
Ego motion, gaze on screen & Preface & \uline{0.656} & \uline{0.612} \\
 & Interleaved & 0.641 & 0.605 \\
\cmidrule(r){1-2} \cmidrule(l){3-4}
Ego motion, vehicle motion & Preface & \uline{0.663} & \uline{0.627} \\
 & Interleaved & 0.623 & 0.544 \\
\bottomrule
\end{tabular}
\end{table}

\begin{figure}[t]
    \centering
    \includegraphics[width=\linewidth]{figures/zeroshot_context_fps.png}
    \caption{Effect of context sampling rate under zero-shot settings (gaze-overlay video input, prefaced context).}
    \label{fig:zeroshot_context_fps}
\end{figure}

\paragraph{Effect of context sampling rate} 
Because the frequency of interleaved context is tied to the video frame rate, we vary the textual sampling rate only for the preface strategy. The results are shown in Fig.~\ref{fig:zeroshot_context_fps}.
The accuracy curves suggest that the optimal sampling rate depends on the specific context type, although denser updates, roughly 4--30~Hz, generally perform better. For macro F1, performance tends to improve monotonically as text density increases. Across sampling rates, \emph{gaze on screen} and \emph{vehicle motion} are among the most beneficial cues and provide relatively consistent gains. One possible explanation is that ego and vehicle states can change rapidly, and higher-frequency updates provide more timely contextual conditioning.



\subsection{Impact of contextual information in fine-tuned settings}
To assess whether the model can learn to exploit additional contextual cues, we fine-tune \qwentwob\ using context-guided QA pairs (Fig.~\ref{fig:context_val_test_comparison}). Incorporating combined context improves performance compared with fine-tuning on egocentric video alone. On the validation set, different context combinations yield comparable performance, with \emph{ego motion + gaze direction} performing slightly better. On the test set, however, the differences become more pronounced. In particular, \emph{ego motion} alone, \emph{ego motion + gaze-on-screen}, and \emph{ego motion + vehicle motion} achieve similar validation accuracy, whereas \emph{ego motion + gaze-on-screen} produces a noticeably larger improvement on the test set. This suggests that gaze signals, when grounded by ego motion, can improve not only predictive accuracy but also generalization.

\begin{figure}
    \centering
    \includegraphics[width=\linewidth]{figures/context_val_test_comparison_xtick.png}
    \caption{Validation vs. test performance of fine-tuned \qwentwob\ models under different context inputs.}
    \label{fig:context_val_test_comparison}
\end{figure}




\subsection{Generalization and Overfitting Check}

\subsubsection{Counterfactual input ablation}
As shown in Tab.~\ref{tab:counterfactual_input_ablation}, the counterfactual ablations demonstrate that the fine-tuned model relies substantially on both visual evidence and temporal ordering. 
Firstly, we checked the model finetuned with only egocentric vision, as shown in Tab.~\ref{tab:counterfactual_egocentric_only}. 
With the original videos, it achieved an accuracy of 0.781 and a macro F1 of 0.778. Removing the video reduced accuracy to 0.483, while replacing it with black frames, random noise, or a mismatched video produced similarly low accuracies of 0.517, 0.485, and 0.525, respectively. This consistent degradation indicates that the model’s performance cannot be explained primarily by textual or class-label priors and that it uses information specific to the corresponding video. Shuffling the frames reduced accuracy to \(0.675 \pm 0.007\), showing that correct temporal ordering contributes meaningfully to prediction while order-independent appearance cues remain informative. Reversing the sequence caused a larger decline to 0.530, suggesting particular sensitivity to motion direction and chronological progression. Overall, these findings indicate that the fine-tuned model combines static visual cues with temporally ordered dynamics rather than relying on a single shortcut.

We also show the results for finetuned model with structured prompt guidance with ego motion and gaze direction, shown in Tab.~\ref{tab:counterfactual_egomotion+gazedirection}. The results show that the model replies on both the visual input as well as the structured ego-motion and gaze-direction context when that context is available. 
With the original video, adding context improves accuracy from 0.707 to 0.830, a gain of 0.123. More importantly, when the frames are replaced with black or random-noise images, the no-context model collapses to the majority-class baseline, with accuracy 0.567 and macro F1 0.362. In contrast, the same corrupted-frame conditions remain much stronger with context, reaching about 0.765 accuracy and 0.763 macro F1. This indicates that the high black/noise performance in the context setting is driven by ego-motion and gaze-direction cues, not by useful visual evidence in the corrupted frames.

At the same time, meaningful visual content still matters. With context, mismatched videos reduce accuracy from 0.830 to 0.671, showing that unrelated visual content actively interferes with prediction rather than being ignored. The temporal ablations suggests that shuffling and reversing frames hurt more without context than with context, suggesting that the ordered ego-motion and gaze sequence partially compensates for disrupted RGB temporal order. Overall, our analysis indicates that the model integrates both visual and structured cues, while relying more heavily on ego-motion and gaze direction when those signals are present.

\begin{table}[t]
\centering
\small 
\caption{Counterfactual input ablation results for the fine-tuned model on the test set. Frame-shuffling results are reported as mean $\pm$ standard deviation over three seeds. Accuracy changes are relative to the original-video condition.}
\label{tab:counterfactual_input_ablation}

\subfloat[Results for the fine-tuned model with egocentric video only, evaluated without structured context at inference time.]{%
\label{tab:counterfactual_egocentric_only}
\begin{tabular}{lccc}
\toprule
\textbf{Input condition} &
\textbf{Acc.} &
\textbf{F1} &
\textbf{$\Delta$Acc.} \\
\midrule
Original video       & \textbf{0.781}     & \textbf{0.778}     & --       \\
% Frame shuffling      & $0.675 \pm 0.007$  & $0.663 \pm 0.007$  & $-0.106$ \\
Frame shuffling      & 0.675              & 0.663              & $-0.106$ \\
Reversed frames      & 0.530              & 0.504              & $-0.251$ \\
Mismatched video     & 0.525              & 0.516              & $-0.257$ \\
Black frames         & 0.517              & 0.515              & $-0.264$ \\
Random-noise frames  & 0.485              & 0.475              & $-0.297$ \\
No video             & 0.483              & 0.481              & $-0.299$ \\
\bottomrule
\end{tabular}}

\vspace{1em}

\subfloat[Results for the fine-tuned model with structured prompt guidance (ego-motion and gaze-direction), evaluated with and without structured context at inference time.]{%
\label{tab:counterfactual_egomotion+gazedirection}
\begin{tabular}{lcccc}
\toprule
% \textbf{Input condition} &
% \textbf{Acc. w/ context} &
% \textbf{F1 w/ context} &
% \textbf{Acc. w/o context} &
% \textbf{F1 w/o context} \\
\textbf{Input condition} &
\textbf{Acc. w/} &
\textbf{F1 w/} &
\textbf{Acc. w/o} &
\textbf{F1 w/o} \\
\midrule
Original video      & \textbf{0.830} & \textbf{0.825} & \textbf{0.707} & \textbf{0.655} \\
% Frame shuffling     & $0.803 \pm 0.004$ & $0.794 \pm 0.004$ & $0.661 \pm 0.001$ & $0.572 \pm 0.002$ \\
Frame shuffling     & 0.803 & 0.794 & 0.661 & 0.572 \\
Reversed frames     & 0.783 & 0.769 & 0.620 & 0.492 \\
Mismatched video    & 0.671 & 0.654 & 0.547 & 0.464 \\
Black frames        & 0.765 & 0.763 & 0.567 & 0.362 \\
Random-noise frames & 0.765 & 0.764 & 0.567 & 0.362 \\
No video            & 0.519 & 0.519 & 0.483 & 0.481 \\
\bottomrule
\end{tabular}}

\end{table}


\subsubsection{Generalization capability retention}

Table~\ref{tab:mmbench_regression} reports MMBench accuracy for the base Qwen3-VL-2B-Instruct model and two LoRA fine-tuned variants, averaged over three random seeds. Both fine-tuned models show a consistent drop in overall accuracy relative to the base model (0.795 base vs. 0.751 and 0.732, a 4.4 and 6.3 percentage-point decrease respectively), indicating that fine-tuning on the crossing-intention task does induce some general-ability forgetting rather than leaving it fully intact. This is consistent with prior findings that single-task LoRA fine-tuning of multimodal LLMs reliably induces measurable forgetting on held-out general benchmarks, even though PEFT methods are known to forget less severely than full-parameter fine-tuning~\cite{zhai2023investigating}. Critically, the magnitude we observe is well within the range the literature treats as a moderate, low-rank fine-tuning regime rather than the catastrophic regime: forgetting has been shown to scale sharply with LoRA rank and training exposure, with severe collapse (e.g., accuracy dropping by tens of points) reported only at high ranks or under excessive fine-tuning, while low-rank, moderately-sized fine-tuning produces comparatively mild degradation~\cite{liang2025lorasculpt}, as in our setup. 

The regression is also not uniform across skill categories: perception-oriented categories (coarse and fine-grained perception) degrade only modestly and consistently across both adapters, while reasoning-oriented categories are more sensitive and more variable — logic reasoning, for instance, drops under one adapter ($-7.0$ points) but is essentially unaffected under the other ($+1.0$ point) — suggesting that the specific fine-tuning configuration, rather than fine-tuning per se, modulates which capabilities are most at risk. 

Taken together, these results indicate that our fine-tuning procedure trades a modest, bounded amount of general visual reasoning ability for substantial gains on the target task, a trade-off consistent with what is reported as typical for low-rank adaptation in the literature, and one that leaves the bulk of the model's general multimodal competence intact.



\begin{table*}[t]
\small
\centering
\caption{MMBench accuracy by L2-category for the base Qwen3-VL-2B-Instruct model versus two LoRA fine-tuned variants. Results are mean $\pm$ std over 3 seeds and 100 samples per category.}
\label{tab:mmbench_regression}
\begin{tabular}{lccc}
\toprule
\multirow{2}{*}{\textbf{Category}} & \multirow{2}{*}{\textbf{Base}} & \textbf{Finetuned} & \textbf{Finetuned} \\
& & \textbf{(egocentric only)} & \textbf{(ego motion+gaze direction)} \\
\midrule
Coarse perception                    & $0.900 \pm 0.046$ & $0.860 \pm 0.027$ & $0.807 \pm 0.012$ \\
Fine-grained perception (instance)   & $0.870 \pm 0.052$ & $0.827 \pm 0.067$ & $0.797 \pm 0.074$ \\
Fine-grained perception (cross-inst.)& $0.730 \pm 0.010$ & $0.707 \pm 0.012$ & $0.670 \pm 0.027$ \\
Attribute reasoning                  & $0.823 \pm 0.042$ & $0.770 \pm 0.044$ & $0.750 \pm 0.020$ \\
Relation reasoning                   & $0.793 \pm 0.015$ & $0.760 \pm 0.030$ & $0.703 \pm 0.050$ \\
Logic reasoning                      & $0.653 \pm 0.075$ & $0.583 \pm 0.067$ & $0.663 \pm 0.031$ \\
\midrule
\textbf{Overall}                     & $\mathbf{0.795 \pm 0.009}$ & $\mathbf{0.751 \pm 0.025}$ & $\mathbf{0.732 \pm 0.015}$ \\
\bottomrule
\end{tabular}
\end{table*}


% \subsection{Discussion and Limitation}

% Our object detection and tracking algorithm might not be powerful enough to accurately and consistently label the same object, which might make the visual prompting underperform the standard prompting. 






\section{Conclusion}
\label{sec:conclusion}

This paper investigated the feasibility of using VLMs to decode pedestrian crossing intention from egocentric videos by formulating intention prediction as a closed-ended VQA task. Our experiments show that, although current off-the-shelf VLMs exhibit non-trivial zero-shot capabilities, they still fall short of a strong task-specific baseline and often struggle with higher-level traffic reasoning. In our setting, neither text-based chain-of-thought prompting nor set-of-marks visual prompting yields performance gains.
In contrast, parameter-efficient fine-tuning remarkably improves performance, with the best fine-tuned VLMs outperforming the CLIP+Transformer baseline by 9\% in accuracy.
Moreover, ego motion emerges as an informative contextual cue in zero-shot settings.
Additional contextual cues (vehicle motion and eye gaze) become more effective once the model is trained to exploit them: in particular, incorporating eye gaze improves both predictive accuracy and generalization.

Our study also has several limitations. First, the dataset is collected in a VR shared-space environment, which may not fully capture the long-tail complexity of real-world urban scenes. 
% (e.g., diverse lighting, sensor noise, unstructured behavior, and rare but safety-critical interactions)
Second, our video input is represented by uniformly sampled frames (often interleaved with timestamps), which introduces a key-frame selection problem and may miss brief but decisive cues. Third, qualitative failures suggest that understanding subtle vehicle intent (e.g., yielding vs. creeping) remains challenging.
Future work should therefore consider (1) traffic-specific multimodal pretraining at scale, (2) adaptive temporal sampling or event-driven frame selection, (3) richer and better-grounded visual representations, and (4) improved evaluation protocols on real-world egocentric data with more rare events and interpretable failure analysis.



\section*{Author contributions}

\noindent
\textbf{Danya Li}: Writing--original draft, Visualization, Validation, Methodology, Investigation, Formal analysis, Data curation, Conceptualization. 
% Wencan Mao: Methodology, Investigation, Formal analysis, Conceptualization. Francisco C. Pereira: Supervision, Resources, Conceptualization. Yu Xiao: Supervision, Resources, Conceptualization. 
\textbf{Xiang Su}: Writing--review \& editing, Supervision, Conceptualization. 
\textbf{Yan Feng}: Writing--review \& editing, Data curation. 
\textbf{Rico Krueger}: Writing--review \& editing, Supervision, Resources, Project administration, Funding acquisition, Conceptualization.



















% \section{PROCEDURE FOR PAPER SUBMISSION}

% \subsection{Selecting a Template (Heading 2)}

% First, confirm that you have the correct template for your paper size. This template has been tailored for output on the US-letter paper size. 
% It may be used for A4 paper size if the paper size setting is suitably modified.

% \subsection{Maintaining the Integrity of the Specifications}

% The template is used to format your paper and style the text. All margins, column widths, line spaces, and text fonts are prescribed; please do not alter them. You may note peculiarities. For example, the head margin in this template measures proportionately more than is customary. This measurement and others are deliberate, using specifications that anticipate your paper as one part of the entire proceedings, and not as an independent document. Please do not revise any of the current designations

% \section{MATH}

% Before you begin to format your paper, first write and save the content as a separate text file. Keep your text and graphic files separate until after the text has been formatted and styled. Do not use hard tabs, and limit use of hard returns to only one return at the end of a paragraph. Do not add any kind of pagination anywhere in the paper. Do not number text heads-the template will do that for you.

% Finally, complete content and organizational editing before formatting. Please take note of the following items when proofreading spelling and grammar:

% \subsection{Abbreviations and Acronyms} Define abbreviations and acronyms the first time they are used in the text, even after they have been defined in the abstract. Abbreviations such as IEEE, SI, MKS, CGS, sc, dc, and rms do not have to be defined. Do not use abbreviations in the title or heads unless they are unavoidable.

% \subsection{Units}

% \begin{itemize}

% \item Use either SI (MKS) or CGS as primary units. (SI units are encouraged.) English units may be used as secondary units (in parentheses). An exception would be the use of English units as identifiers in trade, such as Ò3.5-inch disk driveÓ.
% \item Avoid combining SI and CGS units, such as current in amperes and magnetic field in oersteds. This often leads to confusion because equations do not balance dimensionally. If you must use mixed units, clearly state the units for each quantity that you use in an equation.
% \item Do not mix complete spellings and abbreviations of units: ÒWb/m2Ó or Òwebers per square meterÓ, not Òwebers/m2Ó.  Spell out units when they appear in text: Ò. . . a few henriesÓ, not Ò. . . a few HÓ.
% \item Use a zero before decimal points: Ò0.25Ó, not Ò.25Ó. Use Òcm3Ó, not ÒccÓ. (bullet list)

% \end{itemize}


% \subsection{Equations}

% The equations are an exception to the prescribed specifications of this template. You will need to determine whether or not your equation should be typed using either the Times New Roman or the Symbol font (please no other font). To create multileveled equations, it may be necessary to treat the equation as a graphic and insert it into the text after your paper is styled. Number equations consecutively. Equation numbers, within parentheses, are to position flush right, as in (1), using a right tab stop. To make your equations more compact, you may use the solidus ( / ), the exp function, or appropriate exponents. Italicize Roman symbols for quantities and variables, but not Greek symbols. Use a long dash rather than a hyphen for a minus sign. Punctuate equations with commas or periods when they are part of a sentence, as in

% $$
% \alpha + \beta = \chi \eqno{(1)}
% $$

% Note that the equation is centered using a center tab stop. Be sure that the symbols in your equation have been defined before or immediately following the equation. Use Ò(1)Ó, not ÒEq. (1)Ó or Òequation (1)Ó, except at the beginning of a sentence: ÒEquation (1) is . . .Ó

% \subsection{Some Common Mistakes}
% \begin{itemize}


% \item The word ÒdataÓ is plural, not singular.
% \item The subscript for the permeability of vacuum ?0, and other common scientific constants, is zero with subscript formatting, not a lowercase letter ÒoÓ.
% \item In American English, commas, semi-/colons, periods, question and exclamation marks are located within quotation marks only when a complete thought or name is cited, such as a title or full quotation. When quotation marks are used, instead of a bold or italic typeface, to highlight a word or phrase, punctuation should appear outside of the quotation marks. A parenthetical phrase or statement at the end of a sentence is punctuated outside of the closing parenthesis (like this). (A parenthetical sentence is punctuated within the parentheses.)
% \item A graph within a graph is an ÒinsetÓ, not an ÒinsertÓ. The word alternatively is preferred to the word ÒalternatelyÓ (unless you really mean something that alternates).
% \item Do not use the word ÒessentiallyÓ to mean ÒapproximatelyÓ or ÒeffectivelyÓ.
% \item In your paper title, if the words Òthat usesÓ can accurately replace the word ÒusingÓ, capitalize the ÒuÓ; if not, keep using lower-cased.
% \item Be aware of the different meanings of the homophones ÒaffectÓ and ÒeffectÓ, ÒcomplementÓ and ÒcomplimentÓ, ÒdiscreetÓ and ÒdiscreteÓ, ÒprincipalÓ and ÒprincipleÓ.
% \item Do not confuse ÒimplyÓ and ÒinferÓ.
% \item The prefix ÒnonÓ is not a word; it should be joined to the word it modifies, usually without a hyphen.
% \item There is no period after the ÒetÓ in the Latin abbreviation Òet al.Ó.
% \item The abbreviation Òi.e.Ó means Òthat isÓ, and the abbreviation Òe.g.Ó means Òfor exampleÓ.

% \end{itemize}


% \section{USING THE TEMPLATE}

% Use this sample document as your LaTeX source file to create your document. Save this file as {\bf root.tex}. You have to make sure to use the cls file that came with this distribution. If you use a different style file, you cannot expect to get required margins. Note also that when you are creating your out PDF file, the source file is only part of the equation. {\it Your \TeX\ $\rightarrow$ PDF filter determines the output file size. Even if you make all the specifications to output a letter file in the source - if your filter is set to produce A4, you will only get A4 output. }

% It is impossible to account for all possible situation, one would encounter using \TeX. If you are using multiple \TeX\ files you must make sure that the ``MAIN`` source file is called root.tex - this is particularly important if your conference is using PaperPlaza's built in \TeX\ to PDF conversion tool.

% \subsection{Headings, etc}

% Text heads organize the topics on a relational, hierarchical basis. For example, the paper title is the primary text head because all subsequent material relates and elaborates on this one topic. If there are two or more sub-topics, the next level head (uppercase Roman numerals) should be used and, conversely, if there are not at least two sub-topics, then no subheads should be introduced. Styles named ÒHeading 1Ó, ÒHeading 2Ó, ÒHeading 3Ó, and ÒHeading 4Ó are prescribed.

% \subsection{Figures and Tables}

% Positioning Figures and Tables: Place figures and tables at the top and bottom of columns. Avoid placing them in the middle of columns. Large figures and tables may span across both columns. Figure captions should be below the figures; table heads should appear above the tables. Insert figures and tables after they are cited in the text. Use the abbreviation ÒFig. 1Ó, even at the beginning of a sentence.

% \begin{table}[h]
% \caption{An Example of a Table}
% \label{table_example}
% \begin{center}
% \begin{tabular}{|c||c|}
% \hline
% One & Two\\
% \hline
% Three & Four\\
% \hline
% \end{tabular}
% \end{center}
% \end{table}


%    \begin{figure}[thpb]
%       \centering
%       \framebox{\parbox{3in}{We suggest that you use a text box to insert a graphic (which is ideally a 300 dpi TIFF or EPS file, with all fonts embedded) because, in an document, this method is somewhat more stable than directly inserting a picture.
% }}
%       %\includegraphics[scale=1.0]{figurefile}
%       \caption{Inductance of oscillation winding on amorphous
%        magnetic core versus DC bias magnetic field}
%       \label{figurelabel}
%    \end{figure}
   

% Figure Labels: Use 8 point Times New Roman for Figure labels. Use words rather than symbols or abbreviations when writing Figure axis labels to avoid confusing the reader. As an example, write the quantity ÒMagnetizationÓ, or ÒMagnetization, MÓ, not just ÒMÓ. If including units in the label, present them within parentheses. Do not label axes only with units. In the example, write ÒMagnetization (A/m)Ó or ÒMagnetization {A[m(1)]}Ó, not just ÒA/mÓ. Do not label axes with a ratio of quantities and units. For example, write ÒTemperature (K)Ó, not ÒTemperature/K.Ó

% \section{CONCLUSIONS}

% A conclusion section is not required. Although a conclusion may review the main points of the paper, do not replicate the abstract as the conclusion. A conclusion might elaborate on the importance of the work or suggest applications and extensions. 

% \addtolength{\textheight}{-12cm}   % This command serves to balance the column lengths
                                  % on the last page of the document manually. It shortens
                                  % the textheight of the last page by a suitable amount.
                                  % This command does not take effect until the next page
                                  % so it should come on the page before the last. Make
                                  % sure that you do not shorten the textheight too much.

%%%%%%%%%%%%%%%%%%%%%%%%%%%%%%%%%%%%%%%%%%%%%%%%%%%%%%%%%%%%%%%%%%%%%%%%%%%%%%%%



%%%%%%%%%%%%%%%%%%%%%%%%%%%%%%%%%%%%%%%%%%%%%%%%%%%%%%%%%%%%%%%%%%%%%%%%%%%%%%%%



%%%%%%%%%%%%%%%%%%%%%%%%%%%%%%%%%%%%%%%%%%%%%%%%%%%%%%%%%%%%%%%%%%%%%%%%%%%%%%%%
% \section*{APPENDIX}

% Appendixes should appear before the acknowledgment.

% \section*{ACKNOWLEDGMENT}

% The preferred spelling of the word ÒacknowledgmentÓ in America is without an ÒeÓ after the ÒgÓ. Avoid the stilted expression, ÒOne of us (R. B. G.) thanks . . .Ó  Instead, try ÒR. B. G. thanksÓ. Put sponsor acknowledgments in the unnumbered footnote on the first page.



%%%%%%%%%%%%%%%%%%%%%%%%%%%%%%%%%%%%%%%%%%%%%%%%%%%%%%%%%%%%%%%%%%%%%%%%%%%%%%%%

% References are important to the reader; therefore, each citation must be complete and correct. If at all possible, references should be commonly available publications.

\bibliography{reference}

% \begin{thebibliography}{99}

% \bibitem{c1} G. O. Young, ÒSynthetic structure of industrial plastics (Book style with paper title and editor),Ó 	in Plastics, 2nd ed. vol. 3, J. Peters, Ed.  New York: McGraw-Hill, 1964, pp. 15Ð64.
% \bibitem{c2} W.-K. Chen, Linear Networks and Systems (Book style).	Belmont, CA: Wadsworth, 1993, pp. 123Ð135.
% \bibitem{c3} H. Poor, An Introduction to Signal Detection and Estimation.   New York: Springer-Verlag, 1985, ch. 4.
% \bibitem{c4} B. Smith, ÒAn approach to graphs of linear forms (Unpublished work style),Ó unpublished.
% \bibitem{c5} E. H. Miller, ÒA note on reflector arrays (Periodical styleÑAccepted for publication),Ó IEEE Trans. Antennas Propagat., to be publised.
% \bibitem{c6} J. Wang, ÒFundamentals of erbium-doped fiber amplifiers arrays (Periodical styleÑSubmitted for publication),Ó IEEE J. Quantum Electron., submitted for publication.
% \bibitem{c7} C. J. Kaufman, Rocky Mountain Research Lab., Boulder, CO, private communication, May 1995.
% \bibitem{c8} Y. Yorozu, M. Hirano, K. Oka, and Y. Tagawa, ÒElectron spectroscopy studies on magneto-optical media and plastic substrate interfaces(Translation Journals style),Ó IEEE Transl. J. Magn.Jpn., vol. 2, Aug. 1987, pp. 740Ð741 [Dig. 9th Annu. Conf. Magnetics Japan, 1982, p. 301].
% \bibitem{c9} M. Young, The Techincal Writers Handbook.  Mill Valley, CA: University Science, 1989.
% \bibitem{c10} J. U. Duncombe, ÒInfrared navigationÑPart I: An assessment of feasibility (Periodical style),Ó IEEE Trans. Electron Devices, vol. ED-11, pp. 34Ð39, Jan. 1959.
% \bibitem{c11} S. Chen, B. Mulgrew, and P. M. Grant, ÒA clustering technique for digital communications channel equalization using radial basis function networks,Ó IEEE Trans. Neural Networks, vol. 4, pp. 570Ð578, July 1993.
% \bibitem{c12} R. W. Lucky, ÒAutomatic equalization for digital communication,Ó Bell Syst. Tech. J., vol. 44, no. 4, pp. 547Ð588, Apr. 1965.
% \bibitem{c13} S. P. Bingulac, ÒOn the compatibility of adaptive controllers (Published Conference Proceedings style),Ó in Proc. 4th Annu. Allerton Conf. Circuits and Systems Theory, New York, 1994, pp. 8Ð16.
% \bibitem{c14} G. R. Faulhaber, ÒDesign of service systems with priority reservation,Ó in Conf. Rec. 1995 IEEE Int. Conf. Communications, pp. 3Ð8.
% \bibitem{c15} W. D. Doyle, ÒMagnetization reversal in films with biaxial anisotropy,Ó in 1987 Proc. INTERMAG Conf., pp. 2.2-1Ð2.2-6.
% \bibitem{c16} G. W. Juette and L. E. Zeffanella, ÒRadio noise currents n short sections on bundle conductors (Presented Conference Paper style),Ó presented at the IEEE Summer power Meeting, Dallas, TX, June 22Ð27, 1990, Paper 90 SM 690-0 PWRS.
% \bibitem{c17} J. G. Kreifeldt, ÒAn analysis of surface-detected EMG as an amplitude-modulated noise,Ó presented at the 1989 Int. Conf. Medicine and Biological Engineering, Chicago, IL.
% \bibitem{c18} J. Williams, ÒNarrow-band analyzer (Thesis or Dissertation style),Ó Ph.D. dissertation, Dept. Elect. Eng., Harvard Univ., Cambridge, MA, 1993. 
% \bibitem{c19} N. Kawasaki, ÒParametric study of thermal and chemical nonequilibrium nozzle flow,Ó M.S. thesis, Dept. Electron. Eng., Osaka Univ., Osaka, Japan, 1993.
% \bibitem{c20} J. P. Wilkinson, ÒNonlinear resonant circuit devices (Patent style),Ó U.S. Patent 3 624 12, July 16, 1990. 

% \end{thebibliography}




\end{document}
